% from ega3-1-fr.tex
\oldpage[0\textsubscript{III}]{44}
Avec ces notations :
\begin{proposition}[11.8.3]
\label{0.11.8.3-fr}
\emph{Il existe une homotopie
\[
h:\mathrm C_\bullet(A)\otimes\mathrm C_\bullet(B)
\longrightarrow\mathrm C_\bullet(A)\otimes\mathrm C_\bullet(B),
\]
telle que $h(\sigma,\tau)$ soit combinaison linéaire de couples de simplexes
$(\sigma_i,\tau_i)$ avec $|\sigma_i|\subset|\sigma|$,
$|\tau_i|\subset|\tau|$, et que l'on ait pour $f=j\circ p$,}
\[
\label{0.11.8.3.1-fr}
\tag{11.8.3.1}
f-1=h\circ d+d\circ h.
\]
\end{proposition}

\begin{proof}
Il suffit de définir $h$ sur chaque couple $(\sigma,\tau)$ de simplexes, en
raisonnant par récurrence sur la somme des degrés de $\sigma$ et $\tau$, car
on peut prendre $h=0$ lorsque cette somme est $0$. Soit
\[
\omega=f(\sigma,\tau)-(\sigma,\tau)-h(d(\sigma,\tau));
\]
en vertu de l'hypothèse de récurrence et de la définition de $d$, on a
\[
\omega\in\mathrm C_\bullet(|\sigma|)\otimes\mathrm C_\bullet(|\tau|).
\]
On a
\[
d\omega=f(d(\sigma,\tau))-d(\sigma,\tau)-d(h(d(\sigma,\tau)))
=h(d(d(\sigma,\tau)))=0
\]
en vertu de (\hyperref[0.11.8.3.1-fr]{11.8.3.1}) et de l'hypothèse de
récurrence. Or, on a
$H_q(\mathrm C_\bullet(A))=0$ pour $q>0$ (G, I, 3.7.4), donc aussi
\[
H_q(\mathrm C_\bullet(A)\otimes\mathrm C_\bullet(B))=0\quad(q>0),
\]
en vertu de la formule de Künneth (G, I, 5.5.2). Appliquant cette remarque en
remplaçant $A$ par $|\sigma|$ et $B$ par $|\tau|$, on voit qu'il existe un
élément
\[
\omega'\in\mathrm C_\bullet(|\sigma|)\otimes\mathrm C_\bullet(|\tau|)
\]
tel que $\omega=d\omega'$; prenant $h(\sigma,\tau)=\omega'$, on vérifie
(\hyperref[0.11.8.3.1-fr]{11.8.3.1}) pour le couple $(\sigma,\tau)$ et la
récurrence peut se poursuivre.
\end{proof}

\begin{env}[11.8.4]
\label{0.11.8.4-fr}
Les notations étant celles de (\hyperref[0.11.8.2-fr]{11.8.2}), nous poserons
$(\sigma,\tau)\leq(\sigma',\tau')$ si $|\sigma|\subset|\sigma'|$ et
$|\tau|\subset|\tau'|$. Un \emph{système de coefficients} $\Gamma$ sur
$\Sigma(A)\times\Sigma(B)$ est constitué par une famille
$(\Gamma_{\sigma,\tau})$ de groupes abéliens, où
$\Gamma_{\sigma,\tau}$ ne dépend que des ensembles $|\sigma|$ et $|\tau|$,
et une famille d'homomorphismes
\[
\Gamma_{\sigma,\tau}\longrightarrow\Gamma_{\sigma',\tau'}
\quad\text{pour }(\sigma,\tau)\leq(\sigma',\tau'),
\]
formant un système inductif pour cette relation de préordre. On définit alors
un complexe de cochaînes $\mathrm C^\bullet(A,B;\Gamma)$ comme l'ensemble
des familles $\lambda=(\lambda(\sigma,\tau))$, où $(\sigma,\tau)$ parcourt
$\Sigma(A)\times\Sigma(B)$, avec
$\lambda(\sigma,\tau)\in\Gamma_{\sigma,\tau}$ pour tout couple
$(\sigma,\tau)$. La différentielle est donnée de la façon suivante : si
\[
d(\sigma,\tau)=\sum_i\pm(\sigma_i,\tau_i),
\]
on a $|\sigma_i|\subset|\sigma|$, $|\tau_i|\subset|\tau|$ pour tout $i$, et
l'on prend
\[
d\lambda(\sigma,\tau)=\sum_i\pm\lambda_i(\sigma_i,\tau_i),
\]
où $\lambda_i(\sigma_i,\tau_i)$ désigne l'image canonique de
$\lambda(\sigma_i,\tau_i)$ dans $\Gamma_{\sigma,\tau}$.

Nous dirons qu'une cochaîne $\lambda\in\mathrm C^\bullet(A,B;\Gamma)$ est
\emph{bi-alternée} si $\lambda(\sigma,\tau)=0$ lorsque l'un des deux
simplexes $\sigma$, $\tau$ a deux termes égaux, et si l'on a
\[
\lambda(\pi(\sigma),\tau)=\varepsilon_\pi\lambda(\sigma,\tau),
\qquad
\lambda(\sigma,\pi'(\tau))=\varepsilon_{\pi'}\lambda(\sigma,\tau)
\]
pour des permutations quelconques $\pi$, $\pi'$ des indices. Il est clair
que ces cochaînes engendrent un sous-complexe
$\mathrm L^\bullet(A,B;\Gamma)$ de $\mathrm C^\bullet(A,B;\Gamma)$.
\end{env}

\begin{proposition}[11.8.5]
\label{0.11.8.5-fr}
\emph{L'injection canonique
\[
\mathrm L^\bullet(A,B;\Gamma)\longrightarrow
\mathrm C^\bullet(A,B;\Gamma)
\]
définit un isomorphisme pour la cohomologie de ces deux complexes.}
\end{proposition}

\begin{proof}
Notons que si $p$ et $j$ ont le sens défini dans
(\hyperref[0.11.8.2-fr]{11.8.2}), les applications
\[
{}'p:\lambda\longmapsto\lambda\circ p,\qquad
{}'j:\lambda\longmapsto\lambda\circ j
\]
sont définies dans $\mathrm L^\bullet(A,B;\Gamma)$ et
$\mathrm C^\bullet(A,B;\Gamma)$ respectivement, la première n'étant autre
que l'injection canonique. Comme $\mathop{}'j\circ{}'p$ est l'identité, il
suffit de démontrer que $\mathop{}'p\circ{}'j$ est homotope à l'identité; or,
d'après (\hyperref[0.11.8.3-fr]{11.8.3}),
$\mathop{}'h:\lambda\mapsto\lambda\circ h$ est définie dans
$\mathrm C^\bullet(A,B;\Gamma)$ et on peut donc transposer l'identité
(\hyperref[0.11.8.3.1-fr]{11.8.3.1}), qui fournit le résultat cherché.
\end{proof}

\begin{env}[11.8.6]
\label{0.11.8.6-fr}
